\section{电子探针中的四叶玫瑰线}
\markboth{40. 电子探针中的四叶玫瑰线}{40. 电子探针中的四叶玫瑰线}

1949 年以后，出现了一种分析物质中化学元素成分的精密仪器，叫做“电子探针 X 射线显微分析仪”，简称“电子探针”。只要从样品上取出直径为千分之一毫米的微尘，用一束电子作为“探针”射到微尘上，使它激发出 X 光，然后对这 X 光进行光谱分析，就能知道样品中含有哪些化学元素。

在一些电子探针的仪器内部可以看到两条曲线槽，其中的一条是圆弧槽，另一条就是四叶玫瑰线的一段弧。这两条曲线槽的形状，是根据下面的轨迹问题确定的：

设四边形 $OAMN$ 的一个顶点 $O$ 是定点，边 $ON$、$MN$ 和对角线 $AN$ 都等于定长 $R$，$AM=AO$，顶点 $A$ 在定直线 $OX$ 上滑动，求点 $N$ 的轨迹和 $M$ 的轨迹（参见图 40.1）。

由于 $O$ 为定点，$ON$ 为定长 $R$，立刻知道 $N$ 的轨迹是圆。电子探针中的圆弧槽，就是从这里得来的。

为了求 $M$ 的轨迹，如图 40.1，选取极坐标系，使 $O$ 为极点，$OX$ 为极轴。设 $M$ 的极坐标为 $(\rho, \theta)$，则
\[
  OM = \rho,\quad \angle XOM = \theta.
\]

\noindent
\begin{minipage}[t]{0.44\textwidth}
  \vspace{0pt}
  \centering
  \begin{tikzpicture}[scale=1.12, >=Stealth, line cap=round, line join=round]
    % R = 1.65, theta = 23 deg
    \def\R{1.65}
    \def\th{23}
    \pgfmathsetmacro{\Ax}{2*\R*sin(\th)}
    \pgfmathsetmacro{\Nx}{\R*sin(\th)}
    \pgfmathsetmacro{\Ny}{\R*cos(\th)}
    \pgfmathsetmacro{\Mx}{2*\R*sin(2*\th)*cos(\th)}
    \pgfmathsetmacro{\My}{2*\R*sin(2*\th)*sin(\th)}
    \pgfmathsetmacro{\Dx}{\R*sin(2*\th)*cos(\th)}
    \pgfmathsetmacro{\Dy}{\R*sin(2*\th)*sin(\th)}

    \coordinate (O) at (0,0);
    \coordinate (A) at (\Ax, 0);
    \coordinate (N) at (\Nx, \Ny);
    \coordinate (M) at (\Mx, \My);
    \coordinate (D) at (\Dx, \Dy);

    % 极轴 OX
    \draw (O) -- (2.15, 0);

    % 四叶玫瑰线局部虚线轨迹 rho = 2R sin(2*t)
    \draw[densely dashed, thick, curveblue, domain=0:28, samples=60, variable=\t]
      plot ({2*\R*sin(2*\t)*cos(\t)}, {2*\R*sin(2*\t)*sin(\t)});

    % 四边形与对角线
    \draw[thick] (O) -- (N) -- (M) -- (A) -- cycle;
    \draw (N) -- (A);
    \draw (O) -- (M);

    % 角度 theta
    \draw (0.44, 0) arc (0:\th:0.44);
    \node[font=\scriptsize] at ({0.52*\th}:0.68) {$\theta$};

    \fill (O) circle (1.1pt) node[below left=-1pt, font=\footnotesize] {$O$};
    \fill (A) circle (1.1pt) node[below, font=\footnotesize] {$A$};
    \fill (N) circle (1.1pt) node[above, font=\footnotesize] {$N$};
    \fill (M) circle (1.1pt) node[right, font=\footnotesize] {$M$};
    \fill (D) circle (1.1pt) node[above left=-2pt, font=\footnotesize] {$D$};
    \node[below, font=\footnotesize] at (2.10, 0) {$X$};
  \end{tikzpicture}
  \captionof*{figure}{图\ 40.1}
\end{minipage}
\hfill
\begin{minipage}[t]{0.53\textwidth}
  \vspace{0pt}
  \hspace{2em}由于 $AM=AO$，所以
  \[
    \angle AMO = \angle AOM = \theta.
  \]
  又因为 $NO=NA=NM=R$，所以 $O$、$A$、$M$ 三点都在以 $N$ 为圆心、$R$ 为半径的圆上。在圆 $N$ 中，$\arc{AM}$ 所对的圆周角是 $\angle AOM$，圆心角是 $\angle ANM$，所以
  \[
    \angle ANM = 2\angle AOM = 2\theta.
  \]
\end{minipage}

\medskip
\noindent
设弦 $OM$ 与半径 $NA$ 的交点为 $D$，从直角三角形 $NDM$ 得到
\begin{gather*}
  DM = NM\sin 2\theta = R\sin 2\theta, \\
  \therefore\ OM = 2DM = 2R\sin 2\theta,
\end{gather*}
即
\[
  \rho = 2R\sin 2\theta.
\]
上述方程表明，$M$ 的轨迹是四叶玫瑰线。电子探针里的四叶玫瑰线槽就是从这里来的。

至于在电子探针里为什么要考虑上面的轨迹问题，那就牵涉到物理学中关于光的衍射理论的某些比较专门的知识，不能在这里介绍了。
